\begin{pastexam}[label={px:nchu-111-is-08}]{中興資工所乙組 111 學年度 第 19 題}
（問答）

{[4\%]} Assume that each of the expressions below gives the processing time T(n) spent by an algorithm for
solving a problem of size n. Select the dominant term(s) having the steepest increase in n.
\begin{center}
\begin{tabular}{|c|p{6cm}|}
\hline
Expression & \centering Dominant term(s) \tabularnewline \hline
$5 + 0.001n^3 + 0.025n$ & \\ \hline
$500n + 100n^{0.5} + 50n\log_{10} n$ & \\ \hline
$100n + 0.01n^2$ & \\ \hline
$100n\log_3 n + n^3 + 100n$ & \\ \hline
\end{tabular}
\end{center}
\end{pastexam}

\begin{pastsoln}
主導項是 $n \to \infty$ 時成長最快的那一項，係數大小不影響（資結書第~1.6~節的成長率順序：
$\log n \prec n^{0.5} \prec n \prec n\log n \prec n^2 \prec n^3$）。
\begin{center}
\begin{tabular}{ll}
\toprule
式子 & 主導項 \\
\midrule
$5 + 0.001n^3 + 0.025n$ & $0.001n^3$ \\
$500n + 100n^{0.5} + 50n\log_{10} n$ & $50n\log_{10} n$ \\
$100n + 0.01n^2$ & $0.01n^2$ \\
$100n\log_3 n + n^3 + 100n$ & $n^3$ \\
\bottomrule
\end{tabular}
\end{center}
說明：
\begin{itemize}
\item 第一列：$n^3$ 比 $n$ 和常數成長快，係數 $0.001$ 很小也一樣，$n = 1000$ 時 $0.001n^3 = 10^6$，
$0.025n = 25$。
\item 第二列：$50n\log_{10} n / (500n) = \log_{10} n / 10 \to \infty$，所以 $n\log n$ 項終究超過 $500n$
（$n > 10^{10}$ 之後）。 $n^{0.5}$ 比 $n$ 慢。 對數的底只差常數倍，不影響是不是主導項。
\item 第三列：$0.01n^2 / (100n) = n/10^4 \to \infty$，$n > 10^4$ 之後 $n^2$ 項較大。
\item 第四列：$n^3$ 比 $n\log n$、$n$ 都快。
\end{itemize}
所以四個式子分別是 $\Theta(n^3)$、$\Theta(n\log n)$、$\Theta(n^2)$、$\Theta(n^3)$。
\end{pastsoln}
